% Standalone Conclusion section. Include after the Results (or Discussion)
% via \input{sections/MACE_Conclusion} in the Overleaf main file.

\section{Conclusions}
\label{sec:conclusions}

This study examined how temporally structured control inputs, recursive forecasting, and calibrated uncertainty can be combined to predict multivariate physical-system transients. The Multivariate Adaptive Conformal Ensemble (MACE) framework couples integrated stochastic control-velocity sampling to a ten-member independently recurrent LSTM ensemble and constructs prediction intervals by normalizing forecast errors with ensemble spread. Horizon-specific, per-output complete-trajectory, and joint complete-multivariate trajectory calibrations distinguish the statistical events of interest. These elements were evaluated with a reduced heat-pipe microreactor and balance-of-plant simulator predicting 15 output variables over 200~s.

On independent stochastic and operation-informed test populations, the ensemble generated complete recursive forecasts with pooled normalized mean absolute error and root mean square error of 0.00154 and 0.00359, and 0.00141 and 0.00256, respectively. Under matched stochastic calibration, MACE-Horizon reduced mean normalized full interval width by 8.9\% and interval score by 34.6\% relative to Absolute CP-Horizon; MACE-Trajectory reduced them by 43.9\% and 50.0\% relative to Absolute CP-Trajectory. Both comparisons attained empirical coverage near the nominal 95\% level for their respective objectives, and reductions in width and score were observed for all 15 output variables. Joint complete-multivariate trajectory calibration achieved coverage of 0.9525 with mean normalized width of 0.0169 for MACE, compared with 0.9445 and 0.0398 for joint absolute-error calibration. These results demonstrate improved empirical interval efficiency under the evaluated matched-distribution conditions.

The operation-informed evaluation further established the importance of the calibration distribution. Transferring stochastic thresholds to the same 2{,}000 operation-informed forecasts yielded pooled per-output complete-trajectory coverage of 0.6800 for MACE-Trajectory. Independent operation-informed recalibration raised this value to 0.9549, a paired increase of 0.2749 with a descriptive 95\% whole-trajectory bootstrap interval of [0.2666, 0.2831], without changing the ensemble forecasts. The bootstrap interval describes evaluation-sampling variability conditional on the fixed ensemble and calibration datasets. Recalibration did not uniformly reduce interval width or score. For MACE-Horizon, coverage increased from 0.9479 to 0.9559 while mean normalized width decreased from 0.0073 to 0.0062. For MACE-Trajectory, recovering coverage required width to increase from 0.0121 to 0.0136 and interval score from 0.0123 to 0.0137. Nevertheless, recalibrated MACE-Trajectory retained smaller pooled mean width than Absolute CP-Trajectory (0.0136 versus 0.0246) on the operation-informed distribution.

The observed positive association between ensemble spread and absolute forecast error supports using disagreement as a forecast-dependent interval scale; it is not, by itself, calibrated uncertainty. Standard split-conformal coverage relies on exchangeable calibration and evaluation trajectory samples and does not establish conditional coverage for individual maneuvers or shifted populations. These findings also concern a single reduced simulator, prescribed control trajectories, a fixed 200~s horizon, and fixed trained models and calibration samples rather than operating-reactor measurements. Future investigations should consider higher-fidelity simulators, broader and longer control sequences, variability across new ensembles and calibration samples, and methods for recalibration under changing conditions. The combined results indicate that ensemble disagreement can make conformal intervals more efficient, but reliable coverage requires calibration representative of the trajectories being evaluated.
